% Options for packages loaded elsewhere
\PassOptionsToPackage{unicode}{hyperref}
\PassOptionsToPackage{hyphens}{url}
\documentclass[
  12pt,
]{ctexart}
\usepackage{xcolor}
\usepackage{amsmath,amssymb}
\setcounter{secnumdepth}{-\maxdimen} % remove section numbering
\usepackage{iftex}
\ifPDFTeX
  \usepackage[T1]{fontenc}
  \usepackage[utf8]{inputenc}
  \usepackage{textcomp} % provide euro and other symbols
\else % if luatex or xetex
  \usepackage{unicode-math} % this also loads fontspec
  \defaultfontfeatures{Scale=MatchLowercase}
  \defaultfontfeatures[\rmfamily]{Ligatures=TeX,Scale=1}
\fi
\usepackage{lmodern}
\ifPDFTeX\else
  % xetex/luatex font selection
\fi
% Use upquote if available, for straight quotes in verbatim environments
\IfFileExists{upquote.sty}{\usepackage{upquote}}{}
\IfFileExists{microtype.sty}{% use microtype if available
  \usepackage[]{microtype}
  \UseMicrotypeSet[protrusion]{basicmath} % disable protrusion for tt fonts
}{}
\makeatletter
\@ifundefined{KOMAClassName}{% if non-KOMA class
  \IfFileExists{parskip.sty}{%
    \usepackage{parskip}
  }{% else
    \setlength{\parindent}{0pt}
    \setlength{\parskip}{6pt plus 2pt minus 1pt}}
}{% if KOMA class
  \KOMAoptions{parskip=half}}
\makeatother
\usepackage{longtable,booktabs,array}
\usepackage{calc} % for calculating minipage widths
% Correct order of tables after \paragraph or \subparagraph
\usepackage{etoolbox}
\makeatletter
\patchcmd\longtable{\par}{\if@noskipsec\mbox{}\fi\par}{}{}
\makeatother
% Allow footnotes in longtable head/foot
\IfFileExists{footnotehyper.sty}{\usepackage{footnotehyper}}{\usepackage{footnote}}
\makesavenoteenv{longtable}
\setlength{\emergencystretch}{3em} % prevent overfull lines
\providecommand{\tightlist}{%
  \setlength{\itemsep}{0pt}\setlength{\parskip}{0pt}}
\usepackage{bookmark}
\IfFileExists{xurl.sty}{\usepackage{xurl}}{} % add URL line breaks if available
\urlstyle{same}
\hypersetup{
  hidelinks,
  pdfcreator={LaTeX via pandoc}}

\author{SCX}
\date{2026}

\begin{document}

\section{Spring Algorithm}\label{spring-algorithm}

\begin{quote}
\textbf{Formal name}: Spring Self-Evolving Gatekeeper (SSEG)
\textbf{Abbreviation}: Spring \textbf{Location}:
/g/Xiaogan\_Supercomputing\_data/SCX/theory/self\_evolution/
\textbf{Paper target}: Nature Computational Science (Paper 2 of
two-paper strategy)
\end{quote}

\begin{center}\rule{0.5\linewidth}{0.5pt}\end{center}

\subsection{Why ``Spring''}\label{why-spring}

Spring is the season of \textbf{resurrection}.

\begin{verbatim}
Nature's Spring  →  SCX Spring Algorithm

  Winter dormancy     →  Structures below threshold τ_keep — dormant, not dead
  Spring thaw         →  Gatekeeper S_t updates (new knowledge arrives)
  Re-blooming         →  Old structures re-scored — those crossing τ_keep are RESURRECTED
  Each spring cycle   →  Each iteration: S_t re-evaluates M_t
  Perennial growth    →  Memory bank M_t grows monotonically, never pruned
  Deeper roots        →  DFT/experimental ground truth (physical anchor prevents drift)
\end{verbatim}

\textbf{Tagline}: \emph{``Winter does not kill --- it waits for spring.
Every discarded structure carries the seed of its own resurrection.''}

\begin{center}\rule{0.5\linewidth}{0.5pt}\end{center}

\subsection{Algorithm Sketch}\label{algorithm-sketch}

\begin{verbatim}
Spring(S_0, M_0, D_stream):
  for t = 0, 1, 2, ...:
    // Judge: score incoming data
    for (x, y) in D_t:
      score = S_t(x, y)                    // gatekeeper evaluates
      M_t.insert(x, y, score, metadata)    // store always, never delete
      if score >= τ_keep:
        train_NEP(x, y)                    // only high-score → training
    
    // Update: NEP student learns, gatekeeper evolves
    θ_{t+1} = NEP_update(θ_t, M_t)         // student learns from selected data
    S_{t+1} = Gatekeeper_update(S_t, M_t, f_{θ_{t+1}})  // gatekeeper sees through student's eyes
    
    // Resurrect: re-score old dormant structures
    for (x, y, old_score) in M_t where old_score < τ_keep:
      new_score = S_{t+1}(x, y)
      if new_score >= τ_keep:
        RESURRECT(x, y)                    // spring brings it back to life
        train_NEP(x, y)
\end{verbatim}

\begin{center}\rule{0.5\linewidth}{0.5pt}\end{center}

\subsection{Core Theorems}\label{core-theorems}

\begin{longtable}[]{@{}
  >{\raggedright\arraybackslash}p{(\linewidth - 4\tabcolsep) * \real{0.3214}}
  >{\raggedright\arraybackslash}p{(\linewidth - 4\tabcolsep) * \real{0.3929}}
  >{\raggedright\arraybackslash}p{(\linewidth - 4\tabcolsep) * \real{0.2857}}@{}}
\toprule\noalign{}
\begin{minipage}[b]{\linewidth}\raggedright
Theorem
\end{minipage} & \begin{minipage}[b]{\linewidth}\raggedright
Statement
\end{minipage} & \begin{minipage}[b]{\linewidth}\raggedright
Status
\end{minipage} \\
\midrule\noalign{}
\endhead
\bottomrule\noalign{}
\endlastfoot
\textbf{Spring-1} (Convergence) & Under C1-C7, (S\_t, θ\_t) → (S\emph{,
θ}) a.s. --- the gatekeeper converges to a self-consistent fixed point &
Formalized in \texttt{06\_fixed\_point\_convergence.md} \\
\textbf{Spring-2} (Completeness) & Under physical constraints (finite
data, finite precision, finite compute), ∃ finite T* such that system
reaches ε-approximate fixed point & Formalized in
\texttt{07\_completeness.md} \\
\end{longtable}

\begin{center}\rule{0.5\linewidth}{0.5pt}\end{center}

\subsection{Two-Paper Strategy}\label{two-paper-strategy}

\begin{longtable}[]{@{}
  >{\raggedright\arraybackslash}p{(\linewidth - 4\tabcolsep) * \real{0.3333}}
  >{\raggedright\arraybackslash}p{(\linewidth - 4\tabcolsep) * \real{0.3333}}
  >{\raggedright\arraybackslash}p{(\linewidth - 4\tabcolsep) * \real{0.3333}}@{}}
\toprule\noalign{}
\begin{minipage}[b]{\linewidth}\raggedright
\end{minipage} & \begin{minipage}[b]{\linewidth}\raggedright
Paper 1: SCX
\end{minipage} & \begin{minipage}[b]{\linewidth}\raggedright
Paper 2: Spring
\end{minipage} \\
\midrule\noalign{}
\endhead
\bottomrule\noalign{}
\endlastfoot
\textbf{Title} & \emph{State-Conditioned Expertise: A Complete Theory of
Label Noise Detection} & \emph{Spring: A Self-Evolving Gatekeeper with
Provable Convergence} \\
\textbf{Core} & Thm 1, 2, 3 & Spring-1, Spring-2 \\
\textbf{Target} & JMLR / TMLR & Nature Computational Science \\
\textbf{Story} & ``Noise detection needs state-conditioning --- we prove
necessity + sufficiency'' & ``An AI that learns what `good data' means
--- and brings discarded data back to life'' \\
\end{longtable}

\end{document}
